\documentclass[12pt]{article}
\usepackage{amsmath, amssymb, bm}
\usepackage[margin=2.5cm]{geometry}

\newcommand{\bs}[1]{\boldsymbol{#1}}
\newcommand{\dd}{\,\mathrm{d}}

\title{Physics of the Rod DEM Simulation}
\date{}

\begin{document}
\maketitle

% ──────────────────────────────────────────────────────────
\section{Particle Geometry}

Each particle is a \textit{spherocylinder}: a cylindrical shaft of radius $r$
and shaft length $L_s = L - 2r$ capped by two hemispheres, where $L$ is the
total tip-to-tip length.  The particle orientation is described by a unit vector
$\hat{\bs{e}}$.

\paragraph{Volume.}
\begin{equation}
  V = \tfrac{4}{3}\pi r^3 + \pi r^2 L_s
\end{equation}

\paragraph{Mass.}
\begin{equation}
  m = \rho V
\end{equation}

\paragraph{Moments of inertia} (with $z$ along $\hat{\bs{e}}$):
\begin{align}
  I_{xx} = I_{yy} &= \pi\rho\!\left(
      \tfrac{1}{12}r^2 L_s^3
    + \tfrac{8}{15}r^5
    + \tfrac{1}{3}r^3 L_s^2
    + \tfrac{3}{4}r^4 L_s
  \right) \\
  I_{zz} &= \pi\rho\!\left(
      \tfrac{1}{2}r^4 L_s
    + \tfrac{8}{15}r^5
  \right)
\end{align}

% ──────────────────────────────────────────────────────────
\section{Equations of Motion}

\begin{align}
  m\ddot{\bs{x}} &= \bs{F}_\text{contact} + \bs{F}_\text{lub}
                   + \bs{F}_\text{drag} + \bs{F}_\text{lift} \\
  \bs{I}\dot{\bs{\omega}} &= \bs{\tau}_\text{contact} + \bs{\tau}_\text{lub}
                            + \bs{\tau}_\text{fluid}
\end{align}

where $\bs{I} = \mathrm{diag}(I_{xx}, I_{yy}, I_{zz})$ in the particle frame.

% ──────────────────────────────────────────────────────────
\section{Contact Model (Linear Spring--Dashpot)}

\subsection{Geometry}

Let $d$ be the surface-to-surface separation between the closest points of the
two spherocylinder surfaces and $\hat{\bs{n}}$ the unit normal from particle
$A$ to $B$ at the contact point $\bs{r}_c$.  Contact occurs when the normal
overlap is positive:
\begin{equation}
  \delta_n = r_A + r_B - d > 0
\end{equation}

The relative velocity at the contact point is
\begin{equation}
  \bs{v}_\text{rel} = \bs{v}_B - \bs{v}_A
    + \bs{\omega}_B \times (\bs{r}_c - \bs{r}_{B,\text{eff}})
    - \bs{\omega}_A \times (\bs{r}_c - \bs{r}_{A,\text{eff}})
\end{equation}
decomposed into normal and tangential parts:
\begin{equation}
  v_n = \bs{v}_\text{rel}\cdot\hat{\bs{n}}, \qquad
  \bs{v}_t = \bs{v}_\text{rel} - v_n\hat{\bs{n}}
\end{equation}

\subsection{Effective Parameters}

\begin{equation}
  M_\text{eff} = \frac{m_A m_B}{m_A + m_B}, \qquad
  k_n^\text{eff} = \frac{k_{n,A}\, k_{n,B}}{k_{n,A} + k_{n,B}}, \qquad
  e_n^\text{eff} = \min(e_{n,A},\, e_{n,B})
\end{equation}

\subsection{Contact Time and Damping Coefficients}

\begin{equation}
  t_c = \sqrt{\pi^2 + (\ln e_n^\text{eff})^2}\,
        \sqrt{\frac{M_\text{eff}}{k_n^\text{eff}}}
\end{equation}

\begin{equation}
  \xi_n = -\frac{2 M_\text{eff}}{t_c}\ln e_n^\text{eff}
\end{equation}

The tangential spring constant depends on the contact geometry through the
generalised effective mass $M_t^{-1} = M_\text{eff}^{-1}
+ d_A^2/I_A + d_B^2/I_B$, where $d_{A,B} = |\bs{r}_c - \bs{r}_{A,B}|$:
\begin{equation}
  k_t = \frac{\pi^2 + (\ln e_t^\text{eff})^2}{t_c^2\, M_t^{-1}}, \qquad
  \xi_t = \frac{-2\ln e_t^\text{eff}}{t_c\, M_t^{-1}}
\end{equation}

\subsection{Normal Force}

\begin{equation}
  \bs{F}_n = \bigl(k_n^\text{eff}\,\delta_n + \xi_n\, v_n\bigr)\,(-\hat{\bs{n}})
\end{equation}

\subsection{Tangential Force}

\paragraph{Accumulated tangential displacement.}
The tangential displacement vector $\bs{\delta}_t$ is carried as persistent
per-pair state and updated incrementally each timestep:
\begin{equation}
  \bs{\delta}_t^{n+1} = \bs{\delta}_t^n + \bs{v}_t\,\Delta t
\end{equation}
Two direction vectors are then defined from this updated value:
\begin{equation}
  \hat{\bs{t}}_\delta = \frac{\bs{\delta}_t^{n+1}}{|\bs{\delta}_t^{n+1}|},
  \qquad
  \hat{\bs{t}}_v = \frac{\bs{v}_t}{|\bs{v}_t|}
\end{equation}
with the convention that $\hat{\bs{t}}_v$ falls back to $\hat{\bs{t}}_\delta$
when $|\bs{v}_t| < \varepsilon$, and both are set to $\bs{0}$ when the
respective magnitude is below $\varepsilon = 10^{-8}$.  The two directions
are kept separate because they serve distinct physical roles.

\paragraph{Spring and dashpot forces.}
The spring force is directed along the \emph{accumulated displacement}
$\hat{\bs{t}}_\delta$, since that is the direction in which the spring is
stretched:
\begin{equation}
  \bs{F}_\text{spring} = k_t\,|\bs{\delta}_t^{n+1}|\,\hat{\bs{t}}_\delta
\end{equation}
The viscous damping force is directed along the \emph{instantaneous tangential
velocity}:
\begin{equation}
  \bs{F}_\text{damp} = \xi_t\,\bs{v}_t
\end{equation}
During the particle-generation relaxation phase, the damping force is
suppressed ($\bs{F}_\text{damp} = \bs{0}$) to avoid artificially stiff
behaviour at large overlaps.

\paragraph{Coulomb constraint.}
The pre-limit tangential force is
\begin{equation}
  \bs{F}_t' = \bs{F}_\text{spring} + \bs{F}_\text{damp}
\end{equation}
and the effective friction coefficient is $\mu = (\mu_A + \mu_B)/2$.
The applied tangential force is
\begin{equation}
  \bs{F}_t = \begin{cases}
    \bs{F}_t'
      & |\bs{F}_t'| \le \mu\,|\bs{F}_n| \quad \text{(sticking)} \\[4pt]
    \mu\,|\bs{F}_n|\,\hat{\bs{t}}_v
      & \text{otherwise} \quad \text{(sliding)}
  \end{cases}
\end{equation}
During sliding, the Coulomb force is directed along $\hat{\bs{t}}_v$
(opposing the relative slip velocity), not along $\hat{\bs{t}}_\delta$.

\paragraph{Displacement clamping during sliding.}
When sliding occurs, the stored displacement is immediately clamped so that
the spring is consistent with the Coulomb limit at the next timestep:
\begin{equation}
  \bs{\delta}_t^{n+1} \leftarrow
    \hat{\bs{t}}_v\,\frac{\mu\,|\bs{F}_n|}{k_t}
  \qquad \text{(if sliding)}
\end{equation}
Without this step, $|\bs{\delta}_t|$ grows without bound during sustained
sliding and produces a large spurious force spike when contact returns to
sticking.

\subsection{Torques}

\begin{equation}
  \bs{\tau}_A = (\bs{r}_c - \bs{r}_{A,\text{eff}}) \times (\bs{F}_n + \bs{F}_t),
  \qquad
  \bs{\tau}_B = (\bs{r}_c - \bs{r}_{B,\text{eff}}) \times (-\bs{F}_n - \bs{F}_t)
\end{equation}

% ──────────────────────────────────────────────────────────
\section{Lubrication Forces}

Lubrication is applied for pairs within a surface separation $h \le h_\text{max}(r_A+r_B)$
and is regularised at $h_\text{min}(r_A+r_B)$ to avoid divergence.  Three
geometric configurations contribute, each weighted by inverse separation distance:

\paragraph{End--end (ee).}  Closest hemisphere-to-hemisphere distance $h_{ee}$:
\begin{equation}
  \bs{F}_{ee} = -\eta\,\frac{3\pi L^2}{8\alpha^2\, h_{ee}}\, v_n\,\hat{\bs{n}}
\end{equation}

\paragraph{Shaft--shaft (ss).}  Closest shaft-to-shaft distance $h_{ss}$:
\begin{equation}
  \bs{F}_{ss} = -\eta\,\frac{24\pi L^2}{4\, D_{ss}\, h_{ss}}\, v_n\,\hat{\bs{n}}
\end{equation}
where
\begin{equation}
  D_{ss} = (2\alpha+1)\sqrt{(\alpha^2+\tfrac{1}{4})
            (\hat{\bs{e}}_A\!\times\!\hat{\bs{e}}_B)^2
            + \alpha\bigl(1 + (\hat{\bs{e}}_A\cdot\hat{\bs{e}}_B)^2\bigr)}
\end{equation}
and $\alpha = L/(2r)$ is the aspect ratio.

\paragraph{End--shaft (es).}  Closest hemisphere-to-shaft distance $h_{es}$:
\begin{equation}
  \bs{F}_{es} = -\eta\,\frac{4\pi L^2}{4\sqrt{2}\,\alpha^2\, h_{es}}\, v_n\,\hat{\bs{n}}
\end{equation}

\paragraph{Blending.}  The three contributions are combined via inverse-distance
weighting $w_i = (c_i\, h_i)^{-1}$, $W = \sum_i w_i$:
\begin{equation}
  \bs{F}_\text{lub} = \frac{w_{ee}}{W}\bs{F}_{ee}
                    + \frac{w_{ss}}{W}\bs{F}_{ss}
                    + \frac{w_{es}}{W}\bs{F}_{es}
\end{equation}
where $c_{ee},c_{ss},c_{es}$ are user-supplied manual weights.
The lubrication torque is $\bs{\tau}_\text{lub} = (\bs{r}_c - \bs{r}_\text{eff})\times\bs{F}_\text{lub}$.

% ──────────────────────────────────────────────────────────
\section{Stokes Drag}

The fluid has a linear shear velocity profile $\bs{v}_f(\bs{x}) = \dot{\gamma}z\,\hat{\bs{x}}$.
Drag is computed using the Oberbeck resistance tensor.  In the particle frame ($z$ along $\hat{\bs{e}}$):
\begin{equation}
  K_{xx}^\text{p} = K_{yy}^\text{p} = \frac{16(\beta^2-1)}{(2\beta^2-3)\,\ell(\beta)+\beta},
  \qquad
  K_{zz}^\text{p} = \frac{8(\beta^2-1)}{(2\beta^2-1)\,\ell(\beta)-\beta}
\end{equation}
where $\beta = L/r$ and $\ell(\beta) = \ln(\beta+\sqrt{\beta^2-1})/\sqrt{\beta^2-1}$.

The lab-frame resistance tensor is $\bs{K} = \bs{T}^{-1}\bs{K}^\text{p}\bs{T}$,
where $\bs{T}$ is the transformation matrix from the lab to the particle frame.
The drag force is
\begin{equation}
  \bs{F}_\text{drag} = \eta\pi r\, \mathrm{diag}(K_{xx},K_{yy},K_{zz})\cdot(\bs{v}_f - \bs{v}_p)
\end{equation}
(element-wise product using only the diagonal of $\bs{K}$).

% ──────────────────────────────────────────────────────────
\section{Hydrodynamic Torque}

The fluid torque on a spheroid is (Jeffery / Oberbeck):
\begin{align}
  \tau_x &= \frac{16\pi\eta r^3\beta}{3(\beta_0+\beta^2\gamma_0)}
            \bigl[(1-\beta^2)D_{zy} + (1+\beta^2)(W_{zy}-\omega_x)\bigr] \\
  \tau_y &= \frac{16\pi\eta r^3\beta}{3(\alpha_0+\beta^2\gamma_0)}
            \bigl[(\beta^2-1)D_{xz} + (1+\beta^2)(W_{xz}-\omega_y)\bigr] \\
  \tau_z &= \frac{32\pi\eta r^3\beta}{3(\alpha_0+\beta_0)}
            (W_{yx}-\omega_z)
\end{align}
where $D_{ij}$ and $W_{ij}$ are the symmetric (rate-of-strain) and antisymmetric
(vorticity) parts of the local velocity gradient tensor, and the Oberbeck
coefficients are
\begin{align}
  \alpha_0 = \beta_0 &= \frac{\beta^2}{\beta^2-1}
    + \frac{\beta}{2(\beta^2-1)^{3/2}}\ln\!\frac{\beta-\sqrt{\beta^2-1}}{\beta+\sqrt{\beta^2-1}} \\
  \gamma_0 &= -\frac{2}{\beta^2-1}
    - \frac{\beta}{(\beta^2-1)^{3/2}}\ln\!\frac{\beta-\sqrt{\beta^2-1}}{\beta+\sqrt{\beta^2-1}}
\end{align}

% ──────────────────────────────────────────────────────────
\section{Inertial Lift}

The inertial lift is modelled as
\begin{equation}
  \bs{F}_\text{lift} = c_p\, (\bs{K}\,\bs{L}\,\bs{K})
    \cdot \bs{G}^\text{sgn} \cdot \bs{v}_d
\end{equation}
where $c_p = \pi^2\eta r^2 / \sqrt{\nu}$, $\nu = \eta/\rho_f$,
$\bs{v}_d = \bs{v}_f - \bs{v}_p$, $G^\text{sgn}_{ij} = G_{ij}/\sqrt{|G_{ij}|}$
is a signed-square-root of the velocity gradient, and $\bs{L}$ is an empirical
coefficient matrix
\begin{equation}
  \bs{L} = \begin{pmatrix}
    0.0501 & 0.0329 & 0 \\
    0.0182 & 0.0173 & 0 \\
    0      & 0      & 0.0373
  \end{pmatrix}
\end{equation}

% ──────────────────────────────────────────────────────────
\section{Time Integration (Velocity Verlet)}

\paragraph{Translation.}
\begin{align}
  \bs{a}^{n+1} &= \bs{F}^{n+1}/m \\
  \bs{v}^{n+1} &= \bs{v}^n + \tfrac{1}{2}(\bs{a}^n + \bs{a}^{n+1})\Delta t \\
  \bs{x}^{n+1} &= \bs{x}^n + \bigl(\bs{v}^{n+1} + \tfrac{1}{2}\bs{a}^{n+1}\Delta t\bigr)\Delta t
\end{align}

\paragraph{Rotation.}
\begin{align}
  \bs{\alpha}^{n+1} &= \bs{I}^{-1}\bs{\tau}^{n+1} \\
  \bs{\omega}^{n+1} &= \bs{\omega}^n + \tfrac{1}{2}(\bs{\alpha}^n + \bs{\alpha}^{n+1})\Delta t \\
  \hat{\bs{e}}^{n+1} &= \mathrm{normalize}\!\left[
    \hat{\bs{e}}^n + \bigl(\bs{\omega}^{n+1} + \tfrac{1}{2}\bs{\alpha}^{n+1}\Delta t\bigr)
    \times \hat{\bs{e}}^n\,\Delta t\right]
\end{align}

\paragraph{Lees--Edwards boundary conditions.}
At the $z=0$ and $z=L_z$ boundaries, a Lees--Edwards shear offset $s(t)$ and
velocity $v_\text{LE}$ are applied:
\begin{equation}
  v_x \to v_x \mp v_\text{LE}, \qquad
  x   \to x \mp s(t)
\end{equation}
upon crossing, with periodic wrapping in $x$ and $y$.

% ──────────────────────────────────────────────────────────
\section{Stress Tensor}

The particle-level stress contribution from contact normal ($\bs{F}_n$),
contact tangential ($\bs{F}_t$), and lubrication ($\bs{F}_\text{lub}$) forces
is computed as the Irving--Kirkwood dyadic:
\begin{equation}
  \bs{\sigma}^{(i)} = \bs{F} \otimes \bs{r}_c
\end{equation}
The $xz$ and $zx$ components are accumulated into bin-averaged profiles and
divided by the bin volume to give the bulk stress.

\end{document}
